\documentclass{article}
\usepackage[utf8]{inputenc}  
\usepackage[french]{babel}  
\usepackage{amsmath}  
\usepackage{siunitx}  
\usepackage[version=4]{mhchem}  
\usepackage{tkz-euclide}
\usepackage{pgfplots}

\usetikzlibrary{intersections,angles,quotes}
\usepgfplotslibrary{fillbetween}

\begin{document}

\section{Le cercle trigonométrique}

Le cercle trigonométrique est un outil fondamental pour comprendre les fonctions trigonométriques.  
Il s'agit d'un cercle de rayon 1 centré à l'origine d'un repère orthonormé.

\subsection{Définition}  

Un cercle trigonométrique est un cercle~:
\begin{itemize}
    \item de centre $O$ (l'origine du repère)~;  
    \item de rayon $1$~;  
    \item orienté dans le sens inverse des aiguilles d'une montre (sens trigonométrique).  
\end{itemize}

\subsection{Repérage d'un point sur le cercle}

Tout point $M$ du cercle trigonométrique peut être repéré par un angle $\theta$ (en radians ou en degrés) formé entre l'axe des abscisses et la demi-droite  $[O, M)$.


\subsection{Coordonnées d'un point}  
Pour un point $M$ repéré par un angle $\theta$, ses coordonnées $(x, y)$ sont données par~:  
\[  
x = \cos(\theta) \quad \text{et} \quad y = \sin(\theta)
\]


\end{document}
